\documentclass[10pt,a4paper]{amsart}
\usepackage[margin=19mm]{geometry}
\usepackage{amsmath,amssymb,mathtools}
\usepackage[scr=rsfs,cal=euler]{mathalfa}
\usepackage{xfrac}
\usepackage[unicode,hidelinks,bookmarks=false]{hyperref}
\newcommand{\ie}{i.e.\ }
\newcommand{\cf}{cf.\ }
\newcommand{\resp}{respectively\ }
\newcommand{\Subsection}{\subsection}
\providecommand{\nobreakdash}{\nobreak\hbox{-}}
\setlength{\emergencystretch}{2em}
\multlinegap0pt
\usepackage{amssymb}
\usepackage{dsfont}
\usepackage[shortlabels]{enumitem}
\usepackage{xcolor}
\newtheorem{lemma}{Lemma}
\newcommand{\Id}{\mathrm{Id}}
\newcommand{\R}{\mathbb{R}}
\newcommand*{\dbar}[1]{\bar{\bar{#1}}}
\newcommand*{\tbar}[1]{\bar{\dbar{#1}}}
\title{Full flexibility of isometric immersions of metrics with low H\"older regularity in Poznyak theorem's dimension}
\author{Marta Lewicka}
\date{Source: arXiv:2511.16305v3 (CC BY 4.0)}
\begin{document}
\maketitle

\noindent\emph{Source and licence.} Full flexibility of isometric immersions of metrics with low H\"older regularity in Poznyak theorem's dimension. Version arXiv:2511.16305v3 (CC BY 4.0), distributed by arXiv under the Creative Commons Attribution 4.0 International licence, http://creativecommons.org/licenses/by/4.0/, as recorded in the arXiv licence metadata for this exact version and shown on the abstract page https://arxiv.org/abs/2511.16305v3. Full licence notice: \texttt{licenses/arxiv-2511.16305-CC-BY-4.0.txt}.\par\smallskip

\noindent\textit{Editorial scope.} Counted unit: the article's lemma lem\_propa (source0.tex lines 1569--1624). The article's complete original proof of it is delivered with it (source0.tex lines 1625--1791, one \texttt{proof} environment of 167 lines). No mathematics is written, repaired or re-derived: the statement and the proof are the article's own lines, in the article's own notation, and no retained line is substituted or re-flowed.  Retained uncounted as the source-local context the proof needs: lines 1476--1487; lines 1479--1482; lines 1525--1541.  Nothing was omitted from any retained span, so the package carries no omission record.  No retained line contains a citation, so the wrapper carries no bibliography and no bibliography entry.  No retained line contains a figure environment, an image-inclusion command or drawing code, so no image is delivered and the package declares no asset dependency.  Editorial formatting only, in addition: an AMS \texttt{amsart} preamble that restates the article's own declarations and macros used by the retained lines, namely the environments \texttt{lemma} on separate counters, together with the standard packages those lines need.  The retained environments print in this document as Lemma 1 in order of appearance, whereas the article prints the same items with the numbering of its own full document.  The complete native source of the article as distributed in the arXiv e-print for this exact version is bundled unchanged as \texttt{sources/189509/source0.tex}.
\begin{lemma}\label{lem_det}
Let $u\in\mathcal{C}^1(\bar\omega,\R^4)$, defined on the closure of an
open, bounded set $\omega\subset\R^2$, satisfy:
\begin{equation}\label{immers_gamma}
\frac{1}{\gamma}\Id_2\leq (\nabla u)^T\nabla u \leq \gamma\Id_2
\quad\mbox{ in } \bar\omega,
\end{equation}
 for some $\gamma\geq 1$. Then there holds:
$$\big(\frac{2}{\gamma}\big)^{1/2} \leq |\nabla u| \leq (2\gamma)^{1/2} \quad\mbox{ and }\quad 
\frac{1}{\gamma^2}\leq\det((\nabla u)^T\nabla u)\leq \gamma^2
\quad \mbox{ in } \;\bar\omega.$$
\end{lemma}

\begin{equation}
\frac{1}{\gamma}\Id_2\leq (\nabla u)^T\nabla u \leq \gamma\Id_2
\quad\mbox{ in } \bar\omega,
\end{equation}

\begin{lemma}\label{lem_Tu}
Given $u\in\mathcal{C}^{k+1}(\bar\omega,\R^4)$ posed on the closure of an
open, bounded $\omega\subset\R^2$, with  (\ref{immers_gamma}) for
some $\gamma\geq 1$, consider the tangent frame $T_u=(\nabla u)\big((\nabla u)^T\nabla
u\big)^{-1}\in\mathcal{C}^{k}(\bar\omega,\R^{4\times 2})$. Then:
\begin{itemize}
\item[(i)] $\|T_u\|_0\leq C$ and $\displaystyle{\|\nabla^{(m)}T_u\|_0\leq C
\hspace{-4mm} \sum_{p_1+2p_2+\ldots mp_m=m} \; \prod_{i=1}^m\|\nabla^{(i)}\nabla u\|_0^{p_i}}$
for all $ m=1\ldots k$, where $C$ depends only on $\gamma$ and $k$.
\item[(ii)] If $k\geq 1$ and if, with some constants $\mu>1$, $A<1$ and
  $\bar C>1$, we additionally have:
$$\|\nabla^{(m)}\nabla u\|_0\leq \bar C \mu^{m}A \quad \mbox{ for all } \; m=1\ldots k,$$
then there holds:
$$\|\nabla^{(m)}T_u\|_0\leq C\mu^m A \quad\mbox{ for all } m=1\ldots k,$$
where $C$ depends only on $\gamma$, $\bar C$ and $k$, but not on $\mu, A$.
\end{itemize}  
\end{lemma}

\begin{lemma}\label{lem_propa}
Let $u\in\mathcal{C}^{k+1}(\bar\omega, \R^4)$ be defined on the closure of an
open, bounded set $\omega\subset\R^2$, and satisfy (\ref{immers_gamma}) for
some $\gamma\ge 1$. Let $E_u^1, E_u^2\in\mathcal{C}^{k}(\bar\omega, \R^4)$
be a normal frame of $u$, namely:
\begin{equation}\label{norm_frame}
(\nabla u)^T E_u^1 = (\nabla u)^T E_u^2 = 0,\quad |E^1_u|=|E^2_u|=1.
\quad \langle E^1_u, E^2_u\rangle = 0 \quad\mbox{ in } \bar\omega.
\end{equation}
Then, there exists a constant $\rho_\gamma\in (0,1)$ depending only on $\gamma$, such that
for every $v\in\mathcal{C}^{k+1}(\bar\omega, \R^4)$, satisfying
$\|\nabla v-\nabla u\|_0\leq \rho_\gamma$, there holds:
\begin{itemize}
\item[(i)] $v$ is an immersion and its normal frame $E_v^1,
  E_v^2\in\mathcal{C}^{k}(\bar\omega, \R^4)$, namely: 
$$(\nabla v)^T E_v^1 = (\nabla v)^T E_v^2 = 0 \quad \mbox{and}\quad
|E^1_v|=|E^2_v|=1 \quad\mbox{and} \quad \langle E^1_v, E^2_v\rangle =
0\quad \mbox{in } \bar\omega,$$ 
can be defined via the following formulas:
\begin{equation}\label{propa_def} 
\begin{split}
& E^1_v = \frac{\nu^1_v}{|\nu^1_v|}, \qquad 
E^2_v = \frac{\nu^2_v - \langle \nu_v^2, E^1_v\rangle E^1_v}{|\nu^2_v -
  \langle \nu_v^2, E^1_v\rangle E^1_v|},\\
& \mbox{where } \nu_v^i= \big(\Id_4 - T_v(\nabla v - \nabla u)^T\big)E^i_u
\quad\mbox{for } i=1,2 \\ & \mbox{and} \quad T_v= (\nabla v)((\nabla v)^T\nabla v)^{-1}.
\end{split}
\end{equation}
\item[(ii)] For all $m=0\ldots k$ there holds, with constants $C$ depending only on $\gamma$ and $k$:
\begin{equation*}
\begin{split}    
& \|\nabla^{(m)}(E^1_v-E^1_u)\|_0\\ & \leq C
\hspace{-7mm}\sum_{ p+\sum_{i=1}^mi(p_i+\bar
  p_i + \dbar p_i)=m\hspace{-10mm}} \hspace{-7mm} \|\nabla^{(p)}(\nabla v-\nabla u)\|_0
\prod_{i=1}^m \|\nabla^{(i)}(\nabla v - \nabla u)\|_0^{p_i} \|\nabla^{(i)}\nabla v\|_0^{\bar p_i}
\|\nabla^{(i)}E^1_u\|_0^{\dbar p_i},\\
& \|\nabla^{(m)}(E^2_v-E^2_u)\|_0\\ & \leq C \hspace{-7mm}\sum_{p+\sum_{i=1}^mi(p_i+\bar
  p_i + \dbar p_i+\tbar p_i)=m\hspace{-10mm}} \hspace{-11mm} \|\nabla^{(p)}(\nabla v-\nabla u)\|_0
\prod_{i=1}^m \|\nabla^{(i)}(\nabla v - \nabla u)\|_0^{p_i} \|\nabla^{(i)}\nabla v\|_0^{\bar p_i}
\|\nabla^{(i)}E^1_u\|_0^{\dbar p_i} \|\nabla^{(i)}E^2_u\|_0^{\tbar p_i}.
\end{split}
\end{equation*}
\item[(iii)]  If $k\geq 1$ and if, with some constants $\mu>1, A<1$ and $\bar C>1$, we additionally have:
\begin{equation*}
\begin{split}
& \|\nabla^{(m)}(\nabla v-\nabla u)\|_0\leq \bar C \mu^mA \quad\mbox{ for } m=0\ldots k,
\\ & \|\nabla^{(m)}\nabla u\|_0\leq \bar C\mu^{m}, ~~ \;\;
\|\nabla^{(m)}E^i_u\|_0\leq \bar C\mu^{m}  \quad\mbox{for } m=1\ldots k,\; i=1,2,
\end{split}
\end{equation*}
then there holds:
$$\|\nabla^{(m)}(E^i_v-E^i_u)\|_0\leq C \mu^mA \quad\mbox{ for all }
m=0\ldots k,\; i=1,2,$$
where $C$ depends only on $\gamma, \bar C$ and $k$, but not on $\mu, A$.
\end{itemize}
\end{lemma}

\begin{proof}
{\bf 1.} Observe that, by the first bound in Lemma \ref{lem_det}:
\begin{equation*}
\begin{split}
\|(\nabla v)^T\nabla v - (\nabla u)^T\nabla u\|_0& \leq \|\nabla
v-\nabla u\|_0(\|\nabla u\|_0+\|\nabla v\|_0) \\ & \leq \|\nabla
v-\nabla u\|_0 (2\|\nabla u\|_0+\|\nabla v-\nabla u\|_0) \leq C \rho_\gamma,
\end{split}
\end{equation*}
implying for small $\rho_\gamma$ that $v$ is an immersion and:
\begin{equation}\label{a}
\frac{1}{2\gamma}\Id_2\leq (\nabla v)^T\nabla v\leq 2\gamma \,\Id_2
\quad\mbox{in } \bar\omega.
\end{equation}
Thus, $T_v\in \mathcal{C}^{k}(\bar\omega, \R^{4\times 2})$
is well defined and $\nu_v^i$ are normal to $\nabla v$, as $(\nabla v)^TT_v=\Id_2$ and:
$$(\nabla v)^T\nu_v^i = \big((\nabla v)^T - (\nabla v-\nabla
u)^T\big)E^i_u=(\nabla u)^TE^i_u=0 \quad\mbox{ for } i=1,2.$$
Further, applying Lemma \ref{lem_Tu} to bound $\|T_v\|_0$ by a
constant that only depends on $\gamma$, we get:
\begin{equation}\label{c}
 \|\nu_v^i- E^i_u\|_0\leq  \|T_v\|_0\|\nabla v-\nabla u\|_0\leq C
\rho_\gamma\leq \frac{1}{2} \quad \mbox{ for }\; i=1,2,
\end{equation}
when $\rho_\gamma$ is sufficiently small, so in particular:
$$ |\nu_v^i(x)|\in \big[{1}/{2}, {3}/{2}\big]\quad 
\mbox{ for all } \; x\in\bar\omega, \;  i=1,2$$
implying that $E^1_v$ in (\ref{propa_def}) is well defined. Next, we define and observe:
\begin{equation}\label{jad}
\begin{split}
&\tilde\nu^2_v\doteq \nu_v^2-\langle \nu_v^2, E_v^1\rangle E_v^1\\
& \tilde\nu^2_v- E_u^2 =- T_v(\nabla v-\nabla u)^TE_u^2 -
\langle E_u^2, E_v^1\rangle E_v^1 = - T_v(\nabla v-\nabla u)^TE_u^2 -
\langle E_u^2, E_v^1-E^1_v\rangle E_v^1.
\end{split}
\end{equation}
Since $E_v^1-E_u^1=f(\nu_v^1)- f(E_u^1)$, where the function $f(\xi) = \frac{\xi}{|\xi|}$ is Lipschitz on the
set $\{\xi\in\R^4;\,|\xi|\in [1/2, 3/2]\}$, we consequently obtain, recalling
(\ref{c}) and taking $\rho_\gamma$ small:
\begin{equation*}
\begin{split}
\|\tilde\nu_v^2-E_u^2\|_0&\leq \|T_v\|_0\|\nabla u-\nabla v\|_0 +
\|E_v^1-E_u^1\|_0\leq \|T_v\|_0\|\nabla u-\nabla v\|_0 +
C\|\nu_v^1-E_u^1\|_0 \\ & \leq C \|T_v\|_0\|\nabla u-\nabla v\|_0\leq
C\rho_\gamma\leq \frac{1}{2}.
\end{split}
\end{equation*}
In particular, $E^2_v=f(\tilde \nu^2_v)$ is well defined and we note that:
$$ |\tilde \nu_v^2(x)|\in \big[{1}/{2}, {3}/{2}\big]\quad 
\mbox{ for all } \; x\in\bar\omega.$$
This justifies all the claims in (i), in view of:
$\langle E_v^1, E_v^2\rangle = \frac{1}{|\tilde \nu^2_v|}\big(\langle E_v^1, \nu_v^2\rangle -
  \langle \nu_v^2, E_v^1\rangle |E^1_v|^2\big)=0.$

\smallskip

{\bf 2.} We write:
\begin{equation*}
E^1_v-E^1_u = f(\nu^1_v) - f(E^1_u) = \Big(\int_0^1\nabla
f(t\nu^1_v+(1-t)E^1_u)~\mbox{d}t\Big) (\nu^1_v-E^1_u) 
\end{equation*}
and for each fixed $t\in (0,1)$ and $m=1\ldots k$, we estimate,
using Fa\'a di Bruno's formula:
\begin{equation*}
\|\nabla^{(m)}_x\big((\nabla f)\circ(E_u^1+
t(\nu_v^1-E^1_u))\big)\|_0\leq C \hspace{-4mm}\sum_{p_1+2p_2+\ldots mp_m=m} \;
\prod_{i=1}^m\|\nabla^{(i)}(E_u^1+ t(\nu_v^1-E^1_u))\|_0^{p_i}.
\end{equation*}
The above is trivially valid at $m=0$ as well. Consequently, for all $m=0\ldots k$:
\begin{equation}\label{goa}
\begin{split}
&\|\nabla^{(m)}(E^1_v-E^1_u)\|\leq C \hspace{-6mm}\sum_{p+\sum_{i=1}^mip_i=m}
\hspace{-4mm}\|\nabla^{(p)}(\nu_v^1-E^1_u)\|_0 \prod_{i=1}^m\Big(\|\nabla^{(i)}E_u^1\|_0^{p_i}
+ \|\nabla^{(i)}(\nu_v^1-E^1_u)\|_0^{p_i}\Big)
\\ & \leq C \hspace{-4mm}\sum_{p+\sum_{i=1}^mi(p_i+\dbar p_i)=m}\hspace{-6mm}
\|\nabla^{(p)}(\nu_v^1-E^1_u)\|_0 \prod_{i=1}^m \|\nabla^{(i)}(\nu_v^1-E^1_u)\|_0^{p_i}
\|\nabla^{(i)}E_u^1\|_0^{\dbar p_i}.
\end{split}
\end{equation}
Recall that there holds, for all $m=0\ldots k$: 
\begin{equation}\label{b}
\|\nabla^{(m)}(\nu^i_v-E^i_u)\|_0 \leq C\hspace{-2mm}\sum_{p+q+t=m}\hspace{-2mm}
\|\nabla^{(p)}(\nabla u-\nabla v)\|_0 \|\nabla^{(q)}T_v\|_0  \|\nabla^{(t)}E^i_u\|_0.
\end{equation}
Consequently, (\ref{goa}) yields, for all $m=0\ldots k$:
\begin{equation}\label{goa2}
\begin{split}
&\|\nabla^{(m)}(E^1_v-E^1_u)\|\\ & \leq C \hspace{-4mm}
\sum_{p+\sum_{i=1}^mi(p_i+\bar p_i+\dbar p_i)=m\hspace{-13mm}}\hspace{-5mm}
\|\nabla^{(p)}(\nabla v -\nabla u)\|_0 \prod_{i=1}^m \|\nabla^{(i)}(\nabla v-\nabla u)\|_0^{p_i}
\|\nabla^{(i)}T_v\|_0^{\bar p_i} \|\nabla^{(i)}E_u^1\|_0^{\dbar p_i},
\end{split}
\end{equation}
which implies the first bound in (ii) by Lemma \ref{lem_Tu} (i).
To show the second bound, we separately estimate the two terms
in the right hand side of (\ref{jad}). Firstly, similar to (\ref{b}):
\begin{equation*}
\begin{split}
&\|\nabla^{(m)}(\nu_v^2-E_u^2)\|_0 = \|\nabla^{m}\big(T_v(\nabla v-\nabla u)^TE_u^2\big)\|_0
\\ & \leq C\hspace{-2mm}\sum_{p+q+t=m}\hspace{-2mm} \|\nabla^{(p)}(\nabla v-\nabla u)\|_0
\|\nabla^{(q)}T_v\|_0\|\nabla^{(t)}E^2_u\|_0 \quad\mbox{ for all }\; m=0\ldots k.
\end{split}
\end{equation*}
Secondly, applying (\ref{goa2}) and (\ref{b}), we see that:
\begin{equation*}
\begin{split}
&\|\nabla^{(m)}(\langle E^2_u,E^1_v-E^1_u\rangle E^1_v)\|_0
\leq C\hspace{-4mm}\sum_{p+q+t=m}\hspace{-3mm} \|\nabla^{(p)}(E^1_v-E^1_u)\|_0
\big(\|\nabla^{(q)}E^1_u\|+\|\nabla^{(q)}(E^1_v-E^1_u)\|\big)\|\nabla^{(t)}E^2_u\|_0
\\ & \leq C \hspace{-4mm}
\sum_{p+t+\sum_{i=1}^mi(p_i+\bar p_i+\dbar p_i)=m\hspace{-13mm}}\hspace{-5mm}
\|\nabla^{(p)}(\nabla v -\nabla u)\|_0
\|\nabla^{(t)}E_u^2\|_0\prod_{i=1}^m \|\nabla^{(i)}(\nabla v-\nabla u)\|_0^{p_i}
\|\nabla^{(i)}T_v\|_0^{\bar p_i} \|\nabla^{(i)}E_u^1\|_0^{\dbar p_i}.
\end{split}
\end{equation*}
It follows that the same bound is valid for both terms in
$\tilde\nu_v^2-E_u^2$, so that for  $m=0\ldots k$:
\begin{equation*}
\begin{split}
&\|\nabla^{(m)}(\tilde\nu_v^2-E_u^2)\|_0
\\ & \leq C \hspace{-4mm}
\sum_{p+t+\sum_{i=1}^mi(p_i+\bar p_i+\dbar p_i)=m\hspace{-13mm}}\hspace{-5mm}
\|\nabla^{(p)}(\nabla v -\nabla u)\|_0
\|\nabla^{(t)}E_u^2\|_0\prod_{i=1}^m \|\nabla^{(i)}(\nabla v-\nabla u)\|_0^{p_i}
\|\nabla^{(i)}T_v\|_0^{\bar p_i} \|\nabla^{(i)}E_u^1\|_0^{\dbar p_i}.
\end{split}
\end{equation*}
Writing, as before:
\begin{equation*}
E^2_v-E^2_u = f(\tilde\nu^2_v) - f(E^2_u) = \Big(\int_0^1\nabla
f(t\tilde\nu^2_v+(1-t)E^2_u)~\mbox{d}t\Big) (\tilde\nu^2_v-E^2_u), 
\end{equation*}
and applying the estimate above, we arrive at:
\begin{equation*}
\begin{split}
&\|\nabla^{(m)}(E^2_v-E^2_u)\|\leq C 
\sum_{p+\sum_{i=1}^mip_i=m\hspace{-10mm}}
\|\nabla^{(p)}(\tilde\nu_v^2-E^2_u)\|_0 \prod_{i=1}^m
\big(\|\nabla^{(i)}E_u^2\|_0^{p_i} + \|\nabla^{(i)}(\tilde\nu_v^2-E^2_u)\|_0^{p_i} \big)
\\ & \leq C \hspace{-3mm}
\sum_{p+\sum_{i=1}^mi(p_i+\tbar p_i)=m\hspace{-13mm}}\hspace{-2mm}
\|\nabla^{(p)}(\tilde\nu_v^2-E^2_u)\|_0 \prod_{i=1}^m
\|\nabla^{(i)}(\tilde\nu_v^2-E^2_u)\|_0^{p_i} \|\nabla^{(i)}E_u^2\|_0^{\tbar p_i} 
\\ & \leq C \hspace{-4mm}
\sum_{p+\sum_{i=1}^mi(p_i+\bar p_i+\dbar p_i+\tbar p_i)=m\hspace{-15mm}}\hspace{-5mm}
\|\nabla^{(p)}(\nabla v-\nabla u)\|_0 \prod_{i=1}^m
\|\nabla^{(i)}(\nabla v-\nabla u)\|_0^{p_i} \|\nabla^{(i)}T_v\|_0^{\bar p_i} 
\|\nabla^{(i)}E_u^1\|_0^{\dbar p_i}\|\nabla^{(i)}E_u^2\|_0^{\tbar p_i},
\end{split}
\end{equation*}
which completes the proof of (ii) in virtue of Lemma \ref{lem_Tu} (i).

\smallskip

{\bf 3.} It remains to prove (iii). There, the assumptions and
(ii) yield for $i=1,2$ and $m=0\ldots k$:
\begin{equation*}
\begin{split}
&\|\nabla^{(m)}(E^i_v- E^i_u)\|\leq C \hspace{-4mm}
\sum_{p+\sum_{i=1}^mi(p_i+\bar p_i+\dbar p_i+\tbar p_i)=m\hspace{-5mm}}\hspace{-5mm}
\bar C\mu^p A \; \prod_{i=1}^m (\bar C\mu^i)^{p_i+\bar p_i +\dbar
  p_i}(\bar C \mu^iA)^{\tbar p_i} \leq C \bar C^{m+1}\mu^m A.
\end{split}
\end{equation*}
This ends the proof.
\end{proof}

\end{document}
